\documentclass[sigconf,11pt]{acmart}

\settopmatter{printacmref=false}
\renewcommand\footnotetextcopyrightpermission[1]{}
\setcopyright{none}
\acmConference{ADB Spring 2026}{Spring 2026}{Georgia Tech}
\pagestyle{empty}
\raggedbottom
\widowpenalty=0
\clubpenalty=0
\displaywidowpenalty=0

\usepackage{booktabs}
\usepackage{xurl}
\usepackage{titlesec}
\titlespacing*{\section}{0pt}{1.5ex plus 0.5ex minus 0.5ex}{0.5ex plus 0.2ex}
\titlespacing*{\subsection}{0pt}{1ex plus 0.5ex minus 0.3ex}{0.3ex plus 0.1ex}

\begin{document}

\title{HybridSearch: A SQLite Extension for Hybrid Vector-Keyword Retrieval}

\author{Kanhaiya Madaswar}
\affiliation{\institution{Georgia Institute of Technology}}
\email{}

\maketitle

\noindent\textbf{Repository:} \url{https://github.com/kanhaiya38/TokenSmith}

\vspace{-1em}
\section{Project Overview}

RAG systems perform hybrid retrieval by combining dense vector search (FAISS)
and sparse keyword search (BM25) entirely at the application layer, with no
database engine mediating access. In TokenSmith, a local RAG system for
querying database textbooks, Python loads a FAISS index and a BM25 index as
separate binary files, runs both searches independently, and fuses the results
with Reciprocal Rank Fusion (RRF). This design has two problems. First, there
is no metadata filtering: queries always search the entire corpus, diluting the
top-$k$ budget with irrelevant results even when the user's intent is clearly
scoped to a chapter or section. Second, incremental index updates lack
atomicity: FAISS and BM25 are updated independently, leaving them inconsistent
on failure with no recovery mechanism. Database systems solved both problems via
predicate pushdown and WAL-based index maintenance; RAG systems have not.

I build \textbf{HybridSearch}, a C++ SQLite loadable extension compiled to a
shared library and loaded via \path{sqlite3_load_extension}. The extension
exposes hybrid retrieval as a virtual table, moving keyword search and RRF
fusion into the database engine. Future milestones will move FAISS search fully
into C++ inside the extension using the FAISS C++ API, add metadata predicate
pushdown so the query planner can push \path{chapter = ?} constraints into both
retrievers via \path{xBestIndex} and \path{xFilter}, and add transactional
index updates via SQLite WAL.

\section{Current Progress}

\subsection{C++ SQLite Extension}

The file \path{extension/hybrid_search.cpp} (~290 lines) implements the full
\path{sqlite3_module} virtual table interface. The module registers a virtual
table named \path{hybrid_search} with two visible output columns
(\path{chunk_id}, \path{score}) and three hidden input columns
(\path{faiss_results}, \path{query_text}, \path{top_k}).

\path{xBestIndex} is a minimal stub that marks the three hidden columns as
usable equality constraints. \path{xFilter} is the core of the extension. It
unpacks a caller-supplied BLOB of FAISS-ranked chunk IDs into a
\path{faiss_ranks} map (chunk\_id to 1-based rank), runs an FTS5 \path{MATCH}
query against the \path{chunks_fts} virtual table on the same database
connection to obtain \path{fts5_ranks}, fuses both ranked lists using RRF with
$k=60$ and equal weights of 0.5 each, and materializes up to \path{top_k} rows
into the cursor sorted by descending fused score.

A CMake build file (\path{extension/CMakeLists.txt}) compiles the extension as
\path{hybrid_search.so} inside the conda environment. \path{make build-extension}
drives the full configure and build pipeline.

\subsection{Index Migration}

\path{src/index_migration.py} is a one-time migration script that converts the
existing TokenSmith binary artifacts into a single SQLite database at
\path{index/tokensmith.db}. It reads the chunks pickle (1,825 chunks, 2.7 MB)
and metadata pickle (508 KB), inserts them into a \path{chunks} table with
columns for text, source, section path, and page range, and populates a
\path{chunks_fts} FTS5 virtual table via an FTS5 rebuild command. It then loads
the compiled extension and registers the \path{hybrid_search} virtual table.
The database uses WAL journal mode for concurrent read access.
\path{make migrate-db} runs the migration end-to-end.

\subsection{Python Integration}

A new \path{HybridSQLiteRetriever} class in \path{src/retriever.py} sits
alongside the existing \path{FAISSRetriever} and \path{BM25Retriever}. It
accepts the FAISS index object loaded in Python, runs FAISS search in Python,
packs the resulting ranked \path{chunk_id} list as a little-endian \path{int64}
BLOB via \path{struct.pack}, and passes it to the C++ virtual table together
with the raw query string. The virtual table then handles FTS5 and RRF fusion.
The return format is identical to existing retrievers, a
\path{{chunk_id: score}} dict, making it a drop-in replacement.

A \path{retrieval_backend} flag in \path{config/config.yaml} set to either
\path{"faiss"} or \path{"sqlite"} toggles between the original pipeline and the
new one without breaking existing functionality. The existing benchmark suite
and evaluation metrics remain unchanged.

\section{Challenges and Observations}

\textbf{FAISS and SQLite boundary.} In the current milestone, FAISS search runs
in Python and communicates ranked chunk IDs to the extension via a packed BLOB.
This was a deliberate scoping decision: it allowed the virtual table interface,
FTS5 integration, and RRF fusion to be built and validated end-to-end before
adding FAISS C++ integration. The BLOB serves as the communication channel
between Python's FAISS invocation and C++ RRF computation. The final checkpoint
will eliminate this boundary by moving FAISS search fully inside \path{xFilter}
using the FAISS C++ API, with the serialized index stored as a BLOB in a
\path{faiss_index} table.

\textbf{SQLite extension entry-point naming.} SQLite derives the entry-point
symbol from the \path{.so} filename by inserting underscores between words, so
\path{hybrid_search.so} maps to \path{sqlite3_hybrid_search_init}. An initial
mismatch with \path{sqlite3_hybridsearch_init} caused a \path{dlsym} failure at
load time since the symbol could not be resolved. C++ name mangling was an
additional issue that required wrapping the entry point in
\path{extern "C"}. The Python 3.12 \path{load_extension} \path{entrypoint}
keyword argument was used to make the name explicit and avoid this class of
error going forward.

\textbf{FTS5 versus BM25 scores.} SQLite's FTS5 \path{rank} column uses BM25
internally but returns negative values where more negative means better match,
which differs numerically from the \path{rank_bm25} Python library used by the
existing \path{BM25Retriever}. Since RRF uses only rank positions and not raw
scores, this difference does not affect fusion quality. However, it means FTS5
scores cannot be substituted for BM25 scores in a linear fusion setting, which
is worth noting for any future experiments with weighted linear fusion.

\section{Next Steps}

\textbf{End-to-end benchmarking.} The immediate next step is running the
existing benchmark suite (\path{tests/test_benchmarks.py}) with the SQLite
backend enabled to measure Recall@$k$ and MRR against the FAISS-only baseline
across all 11 benchmark queries. This confirms quality parity and validates that
FTS5 plus FAISS RRF fusion matches the Python-side BM25 plus FAISS baseline. I
will also report the storage footprint of the unified \path{tokensmith.db}
against the current scattered FAISS binary and pickle files.

\textbf{FAISS C++ integration inside the extension.} The central remaining
deliverable is moving FAISS search from Python into C++ \path{xFilter}. This
requires installing \path{libfaiss} with C++ headers via conda-forge,
serializing the FAISS index as a BLOB in a \path{faiss_index} table during
migration, deserializing it on \path{xCreate} and \path{xConnect} into a live
\path{faiss::Index} object, and calling \path{index->search()} directly inside
the extension. This eliminates the BLOB handoff from Python entirely and makes
the extension self-contained: a single SQL query with an embedding BLOB is all
the caller needs to issue.

\textbf{Metadata predicate pushdown.} Once FAISS runs in C++,
\path{xBestIndex} will declare \path{chapter =}, \path{section =}, and page
range as pushable constraints with lower \path{estimatedCost} when filters are
present, following the same contract as R-Tree and FTS5. \path{xFilter} will
translate them into a \path{faiss::IDSelectorBatch} for vector search and a
rowid restriction clause for FTS5. Correctness will be validated by asserting
that filtered results match exhaustive post-retrieval filtering. Latency will be
measured and compared against post-filter at three corpus selectivity levels of
5\%, 20\%, and 50\%.

\textbf{Transactional index updates.} On \path{INSERT} into the corpus, both
\path{chunks_fts} and the serialized FAISS BLOB will be updated within a single
SQLite WAL transaction, giving incremental updates the same atomicity guarantees
as any other SQLite write. I will compare incremental update latency against a
full index rebuild across varying batch sizes.

\clearpage

\appendix
\section*{Appendix: Learning Episode -- Two-Phase Locking}

\textbf{Topic:} Two-Phase Locking (2PL)

\bigskip

\noindent\textbf{Question 1:} What is the growing phase in two-phase locking?

\smallskip
\noindent\textbf{Correct textbook chunk(s):}
\begin{quote}
One protocol that ensures serializability is the two-phase locking protocol.
This protocol requires that each transaction issue lock and unlock requests in
two phases: (1) \textit{Growing phase}: A transaction may obtain locks, but may
not release any lock.
\end{quote}

\bigskip

\noindent\textbf{Question 2:} What is the shrinking phase in two-phase locking,
and what is the lock point?

\smallskip
\noindent\textbf{Correct textbook chunk(s):}
\begin{quote}
\textit{Shrinking phase}: A transaction may release locks, but may not obtain
any new locks. Initially, a transaction is in the growing phase. The transaction
acquires locks as needed. Once the transaction releases a lock, it enters the
shrinking phase, and it can issue no more lock requests. The point in the
schedule where the transaction has obtained its final lock (the end of its
growing phase) is called the lock point of the transaction. Transactions can be
ordered according to their lock points --- this ordering is, in fact, a
serializability ordering for the transactions.
\end{quote}

\bigskip

\noindent\textbf{Question 3:} What constraint does strict 2PL add on top of
basic 2PL, and why?

\smallskip
\noindent\textbf{Correct textbook chunk(s):}
\begin{quote}
Cascading rollbacks can be avoided by a modification of two-phase locking called
the strict two-phase locking protocol. This protocol requires not only that
locking be two phase, but also that all exclusive-mode locks taken by a
transaction be held until that transaction commits. This requirement ensures that
any data written by an uncommitted transaction are locked in exclusive mode until
the transaction commits, preventing any other transaction from reading the data.
\end{quote}

\end{document}
